\section{Phạm vi đề tài}

Đồ án tập trung vào bài toán phát hiện đối tượng bằng YOLO trên hai thiết bị đại
diện là Jetson Nano và Jetson Orin Nano. Phạm vi tối ưu gồm pruning, fine-tuning,
QAT, xuất ONNX và tạo TensorRT engine. Các phép đo tập trung vào accuracy,
latency, FPS, model size, GPU memory, mức sử dụng tài nguyên và công suất theo
khả năng thu thập của thiết bị thực tế.

Ở tầng nền tảng, đồ án xây dựng các chức năng cốt lõi cho Fleet, Model Registry,
Deployment, Command, Runtime, Monitoring, Incident, Knowledge và Copilot.
Hệ thống được đánh giá trong quy mô phòng thí nghiệm với \cbs{số lượng thiết bị
thực tế}, chưa hướng đến thay thế một nền tảng MLOps thương mại ở quy mô hàng
nghìn node.

Đồ án không đặt mục tiêu nghiên cứu một kiến trúc YOLO mới, thiết kế phần cứng
tăng tốc mới hoặc xây dựng mô hình ngôn ngữ nền tảng. Federated learning,
continuous learning, orchestration đa cụm và auto-remediation hoàn toàn tự động
được xem là hướng phát triển sau khi nền tảng cốt lõi đã được kiểm chứng.